% © 2026 Ing. David Jaroš — CC BY-NC-ND 4.0
%
% This work is licensed under a Creative Commons Attribution-NonCommercial-NoDerivatives
% 4.0 International License (CC BY-NC-ND 4.0).
%
% See LICENSE.md for full license text.

\documentclass[12pt,a4paper]{article}
\usepackage[a4paper,margin=2.5cm]{geometry}
\usepackage{amsmath,amssymb,tcolorbox}
\tcbuselibrary{skins,breakable}
\newtcolorbox{statusbox}[1][]{colback=gray!8!white,colframe=black!60,arc=3pt,boxrule=1pt,title=Status Summary,#1}

\title{Canonical Choice for Complex Time in UBT}
\author{Ing.\ David Jaro\v{s}}
\date{2026-05-12}

\begin{document}
\maketitle

\section{The Two Formulations}

\textbf{Version A (with chronofactor).}  In the chronofactor formulation the
fundamental field is written as $\Theta(q,\tau)$ with
\(
\tau = t + i\psi
\).
The archival derivation
\texttt{ARCHIVE/archive\_legacy/ARCHIVE/legacy\_variants/ubt\_compare/DERIVATION\_P2\_complex\_chronofactor\_spectrum.md}
treats $\tau$ as an explicit complex evolution parameter and studies the
spectral consequences of complex-time flow
\(
\partial_\tau\Theta = G\Theta
\).
Physically this makes the phase variable $\psi$ explicit: winding, holonomy,
and phase--entropy mixing are described by an auxiliary circle $S^1_\psi$.
Mathematically this produces an explicit KK/winding tower and permits direct
mode decompositions in \(\psi\).

\medskip

\textbf{Version B (without chronofactor).}  In the chronofactor-free archival
variant the field is written as $\Theta(q)$ only, with all phase data absorbed
into the internal $8$ real degrees of freedom of the biquaternion
\(
\Theta(q)\in\mathbb{H}\otimes\mathbb{C}\cong\mathbb{C}^4
\).
Its README states explicitly that there is no external $\tau$ and that all
phase information is intrinsic to the field.  Physically this removes any
separate imaginary-time fiber; phase observables such as holonomy are treated as
internal channels of $\Theta(q)$.  Mathematically this simplifies the base
manifold to ordinary spacetime but also removes the direct KK interpretation of
winding sectors.

\section{Physical Consequences}

The two formulations differ in four important ways.

\begin{enumerate}
  \item \textbf{Interpretation of $\psi$.}
  Version~A keeps $\psi$ explicit and therefore supports winding-number,
  holonomy, and compact-fiber language.  Version~B treats phase as internal data
  of $\Theta$ and has no separate $\psi$ coordinate.

  \item \textbf{Spectrum.}
  Version~A naturally yields KK towers on $S^1_\psi$, including integer or
  half-integer spectra depending on the boundary condition.  Version~B does not
  provide such a tower directly; three generations must then be justified only
  by internal algebraic structure.

  \item \textbf{Action principle.}
  The canonical action
  \(
  S[\Theta]=\int_{\mathcal{M}} d^4x\,\sqrt{-g}\,\mathcal{L}_{\mathrm{UBT}}
  \)
  is defined on real spacetime \(\mathcal{M}\), so UBT is not written as a
  naive five-dimensional second-time theory.  However, its fully biquaternionic
  form contains
  \(
  \int d^4x\,d\psi
  \)
  with \(\psi\in[0,2\pi R_\psi]\), which means that the canonical formalism
  keeps an explicit compact \(\psi\)-fiber.

  \item \textbf{Canonical bookkeeping of charges and generations.}
  Hypercharge quantisation on the $\psi$-circle and the winding-mode language
  used for generations rely on Version~A notation.  Version~B is conceptually
  cleaner as a four-dimensional field theory, but it weakens the direct
  derivational link to the compact-circle sector used elsewhere in the
  repository.
\end{enumerate}

Therefore the canonical action answers the key interpretive question as follows:
\[
\psi \text{ is best treated as a compact internal fiber coordinate }S^1_\psi,
\]
so the canonical choice is \textbf{not} ``an extra macroscopic second time'' and
\textbf{not} ``purely intrinsic with no coordinate meaning at all.''  It is a
compact KK-like phase coordinate written in the shorthand \(\tau=t+i\psi\).

\section{Canonical Choice}

The canonical formulation is \textbf{Version A}, but with a narrowed
interpretation:
\[
\Theta=\Theta(q,\tau),\qquad \tau=t+i\psi,
\]
where \(\psi\) is a \emph{compact internal phase coordinate} rather than a
second observable time direction.

This verdict follows from the active canonical files:
\begin{itemize}
  \item \texttt{canonical/THEORY/canonical/canonical\_field\_equations.tex}
  writes the master equation in the form \(\Theta(q,\tau)\).
  \item \texttt{canonical/THEORY/canonical/canonical\_action.tex} defines the
  physical action on real spacetime but also gives the biquaternionic form with
  explicit \(\int d\psi\) over the compact fiber.
  \item \texttt{canonical/fields/chronofactor\_operator.tex} defines the winding
  Hilbert space \(L^2(S^1_\psi,\mathbb{C})\), the winding operator, and
  imaginary-time translations, making the \(S^1_\psi\) interpretation explicit.
\end{itemize}

The chronofactor-free version remains valuable as a historical simplification:
it shows how much phase structure can be internalised into \(\Theta\).  But it
is not the active canonical choice because it does not match the audited action
and winding-sector machinery currently used for charge quantisation and
generation bookkeeping.

\begin{statusbox}
Kanonická formulace: \textbf{A --- \(\Theta(q,\tau)\) s \(\tau=t+i\psi\)}\\
Fyzikální zdůvodnění: \textbf{\(\psi\) je kompaktní interní fázová souřadnice \(S^1_\psi\), nikoli druhý makroskopický čas; právě tato kružnice nese winding sektory, holonomii a KK spektrum.}\\
Matematická konzistence: \textbf{[L1]} --- kanonická akce obsahuje explicitní \(\int d\psi\) přes kompaktní fiber a chronofactor operator formalizuje odpovídající Hilbertův prostor.\\
Dopad na GR odvození: \textbf{GR zůstává 4D reálným limitem \(\psi\to0\) / real projection; explicitní \(\psi\)-fiber nerozbíjí Einsteinovo odvození, ale jen rozšiřuje nad ním vnitřní sektor.}\\
Dopad na generace: \textbf{Tři generace se zapisují jako první winding sektory na \(S^1_\psi\); vyšší módy nejsou algebraicky zakázány a musí být interpretovány jako vyšší KK harmonické.}
\end{statusbox}

\end{document}
